\section{Техническое задание}

\subsection{Основание для разработки}

Основанием для разработки является задание на выпускную квалификационную работу бакалавра, посвящённую созданию web-сервиса для размещения объявлений об услугах.

\subsection{Цель и назначение разработки}

Основной целью работы является разработка web-сервиса для автоматизации процессов размещения услуг, управления расписанием исполнителей, поиска подходящих предложений и записи заказчиков на свободные временные интервалы.

Назначение разрабатываемой системы заключается в объединении на одной платформе публикации услуг, отображения доступных слотов, создания заявок на бронирование, ведения переписки по заказу, сопровождения статусов записи, а также накопления отзывов и репутационной информации о пользователях.

Задачами разработки являются:
\begin{itemize}
	\item реализация подсистемы регистрации, авторизации и управления пользовательскими профилями;
	\item разработка каталога услуг с возможностью создания, редактирования и публикации карточек услуг;
	\item реализация механизма управления расписанием исполнителя и автоматического формирования доступных слотов для записи;
	\item разработка подсистемы поиска, фильтрации и бронирования услуг по категориям, городу и доступным датам;
	\item реализация средств взаимодействия между заказчиком и исполнителем, включая обработку заявок и встроенный чат;
	\item создание удобного и понятного web-интерфейса для работы с системой.
\end{itemize}

\subsection{Требования к функциональности}

Программная система должна обеспечивать выполнение следующих функций:
\begin{enumerate}
	\item Регистрация и авторизация пользователей.  
	Система должна обеспечивать создание учётных записей пользователей, вход в систему и завершение сеанса работы. После авторизации пользователю должен предоставляться доступ к функциям в соответствии с его ролью.
	
	\item Работа с профилем пользователя.  
	Система должна позволять пользователю просматривать и редактировать данные своего профиля, включая имя, телефон, город, описание и изображение профиля.
	
	\item Работа с услугами исполнителя.  
	Исполнитель должен иметь возможность создавать, редактировать и скрывать свои услуги. Карточка услуги должна содержать название, описание, категорию, стоимость, длительность, город и изображения.
	
	\item Управление расписанием.  
	Исполнитель должен иметь возможность задавать рабочие интервалы по дням недели, формировать расписание на несколько недель вперёд и отмечать отдельные даты как недоступные.
	
	\item Формирование доступных слотов.  
	На основе расписания система должна автоматически вычислять доступные интервалы времени для записи и отображать пользователю только свободные слоты.
	
	\item Поиск и фильтрация услуг.  
	Система должна обеспечивать поиск услуг по категории, городу и доступным датам. Пользователь должен иметь возможность выбирать несколько дат для поиска услуг с реальной доступностью.
	
	\item Просмотр страницы услуги.  
	Система должна отображать полную информацию об услуге, включая описание, стоимость, длительность, изображения, исполнителя и свободные слоты.
	
	\item Бронирование услуги.  
	Заказчик должен иметь возможность выбрать свободный слот и создать заявку на бронирование услуги с сохранением даты, времени, исполнителя и текущего статуса заявки.
	
	\item Управление заявками.  
	Исполнитель должен иметь возможность просматривать входящие заявки, подтверждать их, отменять, переводить в состояние завершённой услуги или отмечать неявку заказчика.
	
	\item Просмотр собственных бронирований.  
	Заказчик должен иметь возможность просматривать список своих заявок, отслеживать их статусы и переходить к переписке по конкретной записи.
	
	\item Работа с избранным.  
	Заказчик должен иметь возможность добавлять услуги в избранное, удалять их из избранного и просматривать сохранённые предложения в отдельном разделе.
	
	\item Обмен сообщениями.  
	Система должна предоставлять заказчику и исполнителю чат, привязанный к конкретной заявке, с возможностью отправки текстовых сообщений и изображений.
	
	\item Отзывы и репутация.  
	После завершения оказания услуги заказчик и исполнитель должны иметь возможность оставить отзыв друг о друге и оценить результат взаимодействия. На основе накопленных отзывов и оценок в системе должна формироваться репутационная информация о пользователях.
	
	\item Валидация пользовательских данных.  
	Система должна выполнять проверку корректности данных, вводимых при регистрации, авторизации, редактировании профиля, создании услуги, настройке расписания и создании бронирования.
\end{enumerate}

\subsection{Требования к интерфейсу}

Интерфейс системы должен представлять собой web-приложение с единым стилем оформления и понятной навигацией между основными разделами сервиса.

Основные элементы интерфейса включают:
\begin{itemize}
	\item главную страницу с кратким описанием сервиса и переходами к основным разделам;
	\item каталог услуг с карточками предложений и системой фильтрации;
	\item страницу отдельной услуги с описанием, изображениями и свободными слотами записи;
	\item формы регистрации, авторизации и редактирования профиля;
	\item личный кабинет исполнителя для управления услугами, расписанием и входящими заявками;
	\item личный кабинет заказчика для просмотра бронирований и избранных услуг;
	\item интерфейс чата по заявке;
	\item систему уведомлений и сообщений об успешном выполнении действий или возникших ошибках.
\end{itemize}

Интерфейс должен обеспечивать визуальное разделение сценариев работы заказчика и исполнителя. Для исполнителя должны быть доступны функции публикации услуг, настройки расписания и обработки заявок. Для заказчика должны быть доступны функции поиска услуг, выбора временного интервала, создания бронирования и переписки с исполнителем.

Особое внимание должно быть уделено удобству работы с карточками услуг, календарными элементами, заявками и сообщениями. Интерфейс должен быть адаптирован для последовательного выполнения действий без лишних переходов между несвязанными экранами.

На рисунке \ref{mainpage:image} представлен внешний вид главной страницы web-сервиса.

\begin{figure}[H]
	\center{\includegraphics[width=1\linewidth]{main.png}}
	\caption{Главная страница web-сервиса}
	\label{mainpage:image}
\end{figure}

Как показано на рисунке, главная страница содержит краткое описание назначения системы, основные навигационные элементы, а также блоки, позволяющие пользователю быстро перейти к поиску услуг, регистрации или авторизации.

На рисунке \ref{catalog:image} представлен внешний вид страницы каталога услуг.

\begin{figure}[H]
	\center{\includegraphics[width=1\linewidth]{services.png}}
	\caption{Страница каталога услуг}
	\label{catalog:image}
\end{figure}

Страница каталога включает список доступных услуг в виде карточек, панель фильтрации и элементы поиска, позволяющие отбирать предложения по категории, городу и доступным датам.

\subsection{Моделирование вариантов использования}

Для разрабатываемого web-сервиса должна быть построена модель, обеспечивающая наглядное представление вариантов использования системы.

Она помогает в процессе проектирования и детального анализа взаимосвязей между пользователями и функциональными возможностями системы. При построении диаграммы вариантов использования применяется унифицированный язык визуального моделирования UML.

Диаграмма вариантов использования описывает функциональное назначение разрабатываемой системы, то есть показывает, какие действия система будет выполнять в процессе работы. Проектируемая система представляется в виде ряда прецедентов, предоставляемых актёрам, которые взаимодействуют с ней извне. Актёром является внешняя сущность, например пользователь. Прецедент, в свою очередь, служит для описания набора действий, которые система предоставляет пользователю.

В системе предусмотрены три основные роли: неавторизованный пользователь, заказчик и исполнитель. Неавторизованный пользователь может просматривать каталог услуг и страницы отдельных услуг, а также переходить к регистрации и авторизации. Заказчик может искать услуги, просматривать доступные даты, создавать бронирования, переписываться с исполнителем, работать с избранным и оставлять отзыв об исполнителе после завершения услуги. Исполнитель может создавать и редактировать услуги, управлять расписанием, просматривать входящие заявки, изменять их статусы и оставлять отзыв о заказчике после завершения оказания услуги.

На основании анализа предметной области в программе должны быть реализованы следующие прецеденты:
\begin{enumerate}
	\item Вход в систему.
	\item Выход из системы.
	\item Регистрация пользователя.
	\item Просмотр каталога услуг.
	\item Поиск услуг по параметрам и доступным датам.
	\item Просмотр страницы услуги.
	\item Просмотр изображений услуги.
	\item Добавление услуги в избранное.
	\item Просмотр избранных услуг.
	\item Создание бронирования.
	\item Просмотр своих бронирований.
	\item Отмена бронирования заказчиком.
	\item Создание услуги исполнителем.
	\item Редактирование услуги исполнителем.
	\item Загрузка изображений для услуги.
	\item Настройка расписания исполнителя.
	\item Указание исключённых дат.
	\item Просмотр входящих заявок исполнителем.
	\item Подтверждение заявки исполнителем.
	\item Отмена заявки исполнителем.
	\item Отметка услуги как завершённой.
	\item Отметка неявки заказчика.
	\item Переписка по заявке.
	\item Отправка изображения в чат.
	\item Оставление отзыва после завершения услуги.
	\item Редактирование профиля пользователя.
\end{enumerate}

На рисунке \ref{guest-usecase:image} представлена диаграмма вариантов использования для неавторизованного пользователя.

\begin{figure}[H]
	\center{\includegraphics[width=1\linewidth]{nonauth.png}}
	\caption{Диаграмма прецедентов для неавторизованного пользователя}
	\label{guest-usecase:image}
\end{figure}

На рисунке \ref{customer-usecase:image} представлена диаграмма вариантов использования для заказчика.

\begin{figure}[H]
	\center{\includegraphics[width=1\linewidth]{customer.png}}
	\caption{Диаграмма прецедентов для заказчика}
	\label{customer-usecase:image}
\end{figure}

На рисунке \ref{provider-usecase:image} представлена диаграмма вариантов использования для исполнителя.

\begin{figure}[H]
	\center{\includegraphics[width=1\linewidth]{provider.png}}
	\caption{Диаграмма прецедентов для исполнителя}
	\label{provider-usecase:image}
\end{figure}

Представленные диаграммы позволяют наглядно показать различия в функциональных возможностях каждой категории пользователей и зафиксировать основные сценарии взаимодействия с системой.

\subsubsection{Вариант использования «Вход в систему»}

Заинтересованные лица и их требования: зарегистрированный пользователь желает получить доступ к функциям системы в соответствии со своей ролью.

Предусловие: пользователь не аутентифицирован в системе; открыта страница входа.

Постусловие: пользователь аутентифицирован; ему предоставлен доступ к функциям системы в соответствии с его ролью.

Основной успешный сценарий:
\begin{enumerate}
	\item Пользователь открывает страницу входа.
	\item Система отображает форму авторизации.
	\item Пользователь вводит адрес электронной почты и пароль.
	\item Пользователь нажимает кнопку входа.
	\item Система выполняет поиск учётной записи в базе данных.
	\item Система проверяет корректность введённого пароля.
	\item Система создаёт пользовательскую сессию.
	\item Пользователь перенаправляется в личный кабинет.
\end{enumerate}

Альтернативный сценарий:
\begin{enumerate}
	\item На шаге 5 или 6 основного сценария система обнаруживает неверные учётные данные.
	\item Система отображает сообщение об ошибке.
	\item Пользователь остаётся на странице входа и может повторить попытку.
\end{enumerate}

\subsubsection{Вариант использования «Выход из системы»}

Заинтересованные лица и их требования: пользователь желает завершить текущий сеанс работы с системой.

Предусловие: пользователь аутентифицирован в системе.

Постусловие: текущая пользовательская сессия завершена.

Основной успешный сценарий:
\begin{enumerate}
	\item Пользователь нажимает кнопку выхода из системы.
	\item Система завершает текущую пользовательскую сессию.
	\item Система перенаправляет пользователя на главную страницу.
\end{enumerate}

\subsubsection{Вариант использования «Регистрация пользователя»}

Заинтересованные лица и их требования: неавторизованный пользователь желает создать учётную запись для дальнейшей работы с системой.

Предусловие: пользователь не аутентифицирован в системе; открыта страница регистрации.

Постусловие: в системе создана новая учётная запись; пользователь может выполнить вход.

Основной успешный сценарий:
\begin{enumerate}
	\item Пользователь открывает страницу регистрации.
	\item Система отображает форму регистрации.
	\item Пользователь вводит необходимые данные.
	\item Пользователь подтверждает отправку формы.
	\item Система проверяет корректность введённых данных.
	\item Система проверяет отсутствие ранее зарегистрированного пользователя с таким адресом электронной почты.
	\item Система создаёт новую учётную запись.
	\item Система отображает сообщение об успешной регистрации.
	\item Пользователь перенаправляется на страницу входа.
\end{enumerate}

Альтернативный сценарий:
\begin{enumerate}
	\item На шаге 5 основного сценария система обнаруживает ошибки заполнения формы.
	\item Система отображает сообщения об ошибках.
	\item Пользователь исправляет данные и повторно отправляет форму.
\end{enumerate}

\subsubsection{Вариант использования «Просмотр каталога услуг»}

Заинтересованные лица и их требования: пользователь желает ознакомиться со списком услуг, размещённых в системе.

Предусловие: в системе опубликована хотя бы одна активная услуга.

Постусловие: пользователю отображается каталог услуг.

Основной успешный сценарий:
\begin{enumerate}
	\item Пользователь открывает страницу каталога услуг.
	\item Система получает список активных услуг из базы данных.
	\item Система формирует карточки услуг с основной информацией.
	\item Пользователь просматривает отображённый список.
	\item Пользователь при необходимости выбирает одну из услуг для дальнейшего просмотра.
\end{enumerate}

\subsubsection{Вариант использования «Поиск услуг по параметрам и доступным датам»}

Заинтересованные лица и их требования: пользователь желает отфильтровать список услуг и получить только подходящие предложения.

Предусловие: открыта страница каталога услуг.

Постусловие: пользователю отображается отфильтрованный список услуг.

Основной успешный сценарий:
\begin{enumerate}
	\item Пользователь открывает панель фильтрации.
	\item Пользователь выбирает категорию, город и одну или несколько дат.
	\item Пользователь применяет выбранные параметры поиска.
	\item Система обрабатывает введённые критерии.
	\item Система выполняет поиск услуг, по которым имеются свободные интервалы в выбранные даты.
	\item Система отображает обновлённый список услуг.
\end{enumerate}

Альтернативный сценарий:
\begin{enumerate}
	\item На шаге 5 основного сценария система не находит подходящих услуг.
	\item Система отображает сообщение об отсутствии результатов поиска.
\end{enumerate}

\subsubsection{Вариант использования «Просмотр страницы услуги»}

Заинтересованные лица и их требования: пользователь желает ознакомиться с подробной информацией об услуге.

Предусловие: в каталоге доступна выбранная услуга.

Постусловие: открыта страница услуги с полной информацией.

Основной успешный сценарий:
\begin{enumerate}
	\item Пользователь выбирает услугу в каталоге.
	\item Система открывает страницу услуги.
	\item Система отображает название, описание, стоимость, длительность, город, изображения и сведения об исполнителе.
	\item Система отображает доступные слоты записи.
	\item Пользователь принимает решение о дальнейшем действии.
\end{enumerate}

\subsubsection{Вариант использования «Просмотр изображений услуги»}

Заинтересованные лица и их требования: пользователь желает просмотреть изображения услуги в увеличенном виде.

Предусловие: открыта страница услуги; у услуги есть изображения.

Постусловие: изображения отображаются во встроенном просмотрщике.

Основной успешный сценарий:
\begin{enumerate}
	\item Пользователь нажимает на одно из изображений услуги.
	\item Система открывает встроенный просмотрщик поверх страницы.
	\item Система отображает выбранное изображение в полном размере.
	\item Пользователь при необходимости переключается между изображениями.
	\item Пользователь закрывает просмотрщик.
\end{enumerate}

\subsubsection{Вариант использования «Добавление услуги в избранное»}

Заинтересованные лица и их требования: заказчик желает сохранить интересующую услугу для последующего быстрого доступа.

Предусловие: пользователь аутентифицирован как заказчик; открыта страница услуги или карточка услуги.

Постусловие: услуга добавлена в избранное пользователя.

Основной успешный сценарий:
\begin{enumerate}
	\item Заказчик нажимает кнопку добавления в избранное.
	\item Система проверяет, не добавлена ли услуга ранее.
	\item Система сохраняет связь между пользователем и услугой.
	\item Система отображает подтверждение успешного добавления.
\end{enumerate}

\subsubsection{Вариант использования «Просмотр избранных услуг»}

Заинтересованные лица и их требования: заказчик желает просмотреть ранее сохранённые услуги.

Предусловие: пользователь аутентифицирован как заказчик.

Постусловие: пользователю отображается список избранных услуг.

Основной успешный сценарий:
\begin{enumerate}
	\item Заказчик открывает раздел избранного.
	\item Система получает список сохранённых услуг.
	\item Система отображает карточки избранных услуг.
	\item Заказчик при необходимости переходит к нужной услуге или удаляет её из списка.
\end{enumerate}

\subsubsection{Вариант использования «Создание бронирования»}

Заинтересованные лица и их требования: заказчик желает записаться на услугу в свободный временной интервал.

Предусловие: пользователь аутентифицирован как заказчик; по услуге доступны свободные слоты.

Постусловие: в системе создана новая заявка на бронирование.

Основной успешный сценарий:
\begin{enumerate}
	\item Заказчик открывает страницу услуги.
	\item Система отображает свободные слоты записи.
	\item Заказчик выбирает дату и время.
	\item Заказчик подтверждает создание заявки.
	\item Система проверяет актуальность выбранного слота.
	\item Система сохраняет заявку в базе данных.
	\item Система присваивает заявке статус ожидания.
	\item Система отображает сообщение об успешном создании заявки.
\end{enumerate}

Альтернативный сценарий:
\begin{enumerate}
	\item На шаге 5 основного сценария система определяет, что выбранный слот уже занят.
	\item Система отображает сообщение о невозможности бронирования выбранного интервала.
	\item Заказчик возвращается к выбору другого слота.
\end{enumerate}

\subsubsection{Вариант использования «Просмотр своих бронирований»}

Заинтересованные лица и их требования: заказчик желает просматривать созданные им заявки и их текущие статусы.

Предусловие: пользователь аутентифицирован как заказчик; у него существует хотя бы одно бронирование.

Постусловие: пользователю отображается список его заявок.

Основной успешный сценарий:
\begin{enumerate}
	\item Заказчик открывает раздел своих бронирований.
	\item Система получает заявки текущего пользователя.
	\item Система отображает дату, время, услугу, исполнителя и статус по каждой заявке.
	\item Заказчик при необходимости открывает детальный просмотр выбранной заявки.
\end{enumerate}

\subsubsection{Вариант использования «Отмена бронирования заказчиком»}

Заинтересованные лица и их требования: заказчик желает отменить ранее созданную заявку.

Предусловие: пользователь аутентифицирован как заказчик; заявка ещё не завершена.

Постусловие: статус заявки изменён на отменённый.

Основной успешный сценарий:
\begin{enumerate}
	\item Заказчик открывает выбранную заявку.
	\item Заказчик выбирает действие отмены бронирования.
	\item Система запрашивает подтверждение действия.
	\item Заказчик подтверждает отмену.
	\item Система обновляет статус заявки.
	\item Система отображает актуальное состояние бронирования.
\end{enumerate}

\subsubsection{Вариант использования «Создание услуги исполнителем»}

Заинтересованные лица и их требования: исполнитель желает разместить новую услугу в системе.

Предусловие: пользователь аутентифицирован как исполнитель.

Постусловие: создана новая карточка услуги.

Основной успешный сценарий:
\begin{enumerate}
	\item Исполнитель открывает раздел управления услугами.
	\item Исполнитель выбирает действие создания новой услуги.
	\item Система отображает форму создания услуги.
	\item Исполнитель заполняет необходимые поля.
	\item Исполнитель сохраняет введённые данные.
	\item Система проверяет корректность заполнения формы.
	\item Система создаёт карточку услуги.
	\item Новая услуга отображается в списке услуг исполнителя.
\end{enumerate}

\subsubsection{Вариант использования «Редактирование услуги исполнителем»}

Заинтересованные лица и их требования: исполнитель желает изменить сведения о своей услуге.

Предусловие: пользователь аутентифицирован как исполнитель; у него имеется хотя бы одна услуга.

Постусловие: изменения по услуге сохранены в системе.

Основной успешный сценарий:
\begin{enumerate}
	\item Исполнитель открывает список своих услуг.
	\item Исполнитель выбирает нужную услугу.
	\item Система отображает форму редактирования.
	\item Исполнитель изменяет необходимые параметры.
	\item Исполнитель сохраняет изменения.
	\item Система проверяет корректность введённых данных.
	\item Система обновляет сведения об услуге.
\end{enumerate}

\subsubsection{Вариант использования «Загрузка изображений для услуги»}

Заинтересованные лица и их требования: исполнитель желает добавить изображения к карточке услуги.

Предусловие: пользователь аутентифицирован как исполнитель; услуга уже создана.

Постусловие: изображения прикреплены к карточке услуги.

Основной успешный сценарий:
\begin{enumerate}
	\item Исполнитель открывает форму редактирования услуги.
	\item Исполнитель выбирает изображения для загрузки.
	\item Система проверяет формат и размер файлов.
	\item Система загружает изображения на сервер.
	\item Система связывает изображения с карточкой услуги.
	\item Изображения становятся доступными на странице услуги.
\end{enumerate}

\subsubsection{Вариант использования «Настройка расписания исполнителя»}

Заинтересованные лица и их требования: исполнитель желает определить рабочие интервалы для оказания услуги.

Предусловие: пользователь аутентифицирован как исполнитель; у него имеется опубликованная услуга.

Постусловие: расписание сохранено; система формирует доступные слоты записи.

Основной успешный сценарий:
\begin{enumerate}
	\item Исполнитель открывает раздел управления расписанием.
	\item Система отображает форму задания рабочих дней и временных интервалов.
	\item Исполнитель указывает доступные периоды работы.
	\item Исполнитель сохраняет изменения.
	\item Система проверяет корректность интервалов.
	\item Система сохраняет расписание.
	\item Система формирует свободные слоты записи.
\end{enumerate}

\subsubsection{Вариант использования «Указание исключённых дат»}

Заинтересованные лица и их требования: исполнитель желает исключить отдельные даты из расписания.

Предусловие: пользователь аутентифицирован как исполнитель; расписание уже задано.

Постусловие: выбранные даты исключены из числа доступных для записи.

Основной успешный сценарий:
\begin{enumerate}
	\item Исполнитель открывает календарь исключённых дат.
	\item Исполнитель выбирает одну или несколько дат.
	\item Исполнитель подтверждает сохранение изменений.
	\item Система сохраняет список недоступных дат.
	\item Система обновляет доступные слоты записи.
\end{enumerate}

\subsubsection{Вариант использования «Просмотр входящих заявок исполнителем»}

Заинтересованные лица и их требования: исполнитель желает просматривать поступившие заявки по своим услугам.

Предусловие: пользователь аутентифицирован как исполнитель; по его услугам имеются заявки.

Постусловие: исполнителю отображается список заявок.

Основной успешный сценарий:
\begin{enumerate}
	\item Исполнитель открывает раздел заявок.
	\item Система получает бронирования по услугам исполнителя.
	\item Система отображает список заявок с основной информацией.
	\item Исполнитель выбирает интересующую заявку для детального просмотра.
\end{enumerate}

\subsubsection{Вариант использования «Подтверждение заявки исполнителем»}

Заинтересованные лица и их требования: исполнитель желает подтвердить возможность оказания услуги по заявке.

Предусловие: пользователь аутентифицирован как исполнитель; заявка находится в статусе ожидания.

Постусловие: заявка переведена в статус подтверждённой.

Основной успешный сценарий:
\begin{enumerate}
	\item Исполнитель открывает заявку.
	\item Исполнитель выбирает действие подтверждения.
	\item Система запрашивает подтверждение действия.
	\item Исполнитель подтверждает выполнение операции.
	\item Система обновляет статус заявки.
\end{enumerate}

\subsubsection{Вариант использования «Отмена заявки исполнителем»}

Заинтересованные лица и их требования: исполнитель желает отменить заявку, если оказание услуги невозможно.

Предусловие: пользователь аутентифицирован как исполнитель; заявка ещё не завершена.

Постусловие: заявка переведена в статус отменённой.

Основной успешный сценарий:
\begin{enumerate}
	\item Исполнитель открывает заявку.
	\item Исполнитель выбирает действие отмены заявки.
	\item Система запрашивает подтверждение действия.
	\item Исполнитель подтверждает отмену.
	\item Система изменяет статус заявки.
\end{enumerate}

\subsubsection{Вариант использования «Отметка услуги как завершённой»}

Заинтересованные лица и их требования: исполнитель желает зафиксировать факт успешного оказания услуги.

Предусловие: пользователь аутентифицирован как исполнитель; заявка подтверждена.

Постусловие: заявка переведена в статус завершённой услуги.

Основной успешный сценарий:
\begin{enumerate}
	\item Исполнитель открывает подтверждённую заявку.
	\item Исполнитель выбирает действие завершения услуги.
	\item Система запрашивает подтверждение действия.
	\item Исполнитель подтверждает завершение.
	\item Система обновляет статус заявки.
\end{enumerate}

\subsubsection{Вариант использования «Отметка неявки заказчика»}

Заинтересованные лица и их требования: исполнитель желает зафиксировать отсутствие заказчика в согласованное время.

Предусловие: пользователь аутентифицирован как исполнитель; заявка подтверждена, но услуга не была оказана из-за отсутствия клиента.

Постусловие: заявка переведена в статус неявки.

Основной успешный сценарий:
\begin{enumerate}
	\item Исполнитель открывает заявку.
	\item Исполнитель выбирает действие отметки неявки.
	\item Система запрашивает подтверждение действия.
	\item Исполнитель подтверждает выполнение операции.
	\item Система сохраняет новый статус заявки.
\end{enumerate}

\subsubsection{Вариант использования «Переписка по заявке»}

Заинтересованные лица и их требования: заказчик и исполнитель желают обмениваться сообщениями по конкретной заявке.

Предусловие: пользователь аутентифицирован; у пользователя есть доступ к выбранной заявке.

Постусловие: новое сообщение сохранено и отображено в истории чата.

Основной успешный сценарий:
\begin{enumerate}
	\item Пользователь открывает чат по заявке.
	\item Система загружает историю сообщений.
	\item Пользователь вводит текст сообщения.
	\item Пользователь отправляет сообщение.
	\item Система сохраняет сообщение.
	\item Система отображает новое сообщение в истории переписки.
\end{enumerate}

\subsubsection{Вариант использования «Отправка изображения в чат»}

Заинтересованные лица и их требования: пользователь желает отправить изображение в рамках переписки по заявке.

Предусловие: пользователь аутентифицирован; открыт чат по заявке.

Постусловие: изображение загружено и отображено в истории чата.

Основной успешный сценарий:
\begin{enumerate}
	\item Пользователь нажимает кнопку добавления изображения.
	\item Пользователь выбирает файл на устройстве.
	\item Система проверяет допустимость формата и размера файла.
	\item Система загружает изображение на сервер.
	\item Система создаёт сообщение с вложением.
	\item Изображение отображается в истории переписки.
\end{enumerate}

\subsubsection{Вариант использования «Оставление отзыва после завершения услуги»}

Заинтересованные лица и их требования: заказчик или исполнитель желает оставить отзыв о второй стороне после завершения оказания услуги и зафиксировать результат взаимодействия.

Предусловие: пользователь аутентифицирован в системе; заявка переведена в статус завершённой услуги; пользователь является участником данной заявки.

Постусловие: отзыв сохранён в системе; репутационные сведения пользователя, на которого оставлен отзыв, обновлены.

Основной успешный сценарий:
\begin{enumerate}
	\item Пользователь открывает завершённую заявку.
	\item Система отображает форму оставления отзыва.
	\item Пользователь указывает оценку и текст комментария.
	\item Пользователь подтверждает сохранение отзыва.
	\item Система проверяет корректность введённых данных.
	\item Система сохраняет отзыв в базе данных.
	\item Система обновляет репутационные показатели пользователя, на которого был оставлен отзыв.
\end{enumerate}

Альтернативный сценарий:
\begin{enumerate}
	\item На шаге 5 система определяет, что текущий пользователь уже оставил отзыв по данной заявке.
	\item Система отображает сообщение о невозможности повторного добавления отзыва.
	\item Новый отзыв не создаётся.
\end{enumerate}

\subsubsection{Вариант использования «Редактирование профиля пользователя»}

Заинтересованные лица и их требования: пользователь желает изменить сведения своего профиля.

Предусловие: пользователь аутентифицирован в системе.

Постусловие: изменения профиля сохранены в базе данных.

Основной успешный сценарий:
\begin{enumerate}
	\item Пользователь открывает страницу редактирования профиля.
	\item Система отображает текущие сведения.
	\item Пользователь изменяет доступные поля.
	\item Пользователь сохраняет внесённые изменения.
	\item Система проверяет корректность введённых данных.
	\item Система обновляет профиль пользователя.
	\item Система отображает сообщение об успешном сохранении изменений.
\end{enumerate}

\subsection{Требования к надёжности}

Система должна обеспечивать надёжное хранение и обработку данных при различных режимах работы.

Целостность данных.  
Хранение данных в системе должно осуществляться с использованием системы управления базами данных MySQL. При выполнении операций создания, изменения и удаления записей должна обеспечиваться корректность обработки данных, исключающая появление противоречивых, дублирующихся или частично сохранённых записей. Для критически важных операций, связанных с бронированиями, отзывами и изменением статусов заявок, должны использоваться механизмы, обеспечивающие согласованность данных на уровне базы данных и серверной логики приложения.

Устойчивость к ошибкам обработки.  
При возникновении ошибок в процессе выполнения пользовательских действий система не должна повреждать уже сохранённые данные. В случае некорректного ввода, сбоя при загрузке файлов, ошибки при выполнении запроса к базе данных или иной нештатной ситуации пользователю должно выводиться соответствующее сообщение, а операция должна завершаться корректно без нарушения общей работоспособности системы.

Контроль целостности на уровне приложения.  
На уровне web-приложения должна быть реализована валидация вводимых данных. Проверке подлежат регистрационные данные пользователя, параметры профиля, сведения об услугах, стоимость, длительность, даты, временные интервалы, содержимое сообщений и иные поля, вводимые через интерфейс системы. Это позволяет предотвратить попадание некорректных данных в базу.

Устойчивость к пользовательским ошибкам.  
Система должна обеспечивать устойчивость к ошибочным действиям пользователя, связанным с вводом некорректных значений, попытками повторной отправки формы, выбором уже занятого слота или загрузкой недопустимых файлов. Во всех подобных случаях система должна корректно обрабатывать ситуацию и сохранять работоспособность.

Логирование ошибок.  
При возникновении критических ошибок и исключительных ситуаций система должна обеспечивать возможность их журналирования на стороне сервера. Записи журнала должны содержать сведения, достаточные для диагностики причины сбоя и последующего анализа работы приложения.

Резервное копирование.  
Для обеспечения сохранности информации должна быть предусмотрена возможность резервного копирования базы данных MySQL и пользовательских файлов, загружаемых в систему. Восстановление данных должно осуществляться с использованием резервных копий базы данных и файловой структуры проекта.

\subsection{Требования к информационной безопасности}

Система должна обеспечивать защиту данных от несанкционированного доступа и разграничение прав пользователей.

Авторизация.  
Доступ к функциям, связанным с редактированием профиля, созданием услуг, настройкой расписания, бронированием, отправкой сообщений и оставлением отзывов, должен предоставляться только после прохождения процедуры авторизации. Для неавторизованных пользователей должны быть доступны только общедоступные разделы, такие как главная страница, каталог услуг и страницы отдельных услуг.

Хранение паролей.  
Пароли пользователей не должны храниться в базе данных в открытом виде. Для их хранения должен использоваться механизм хеширования. При авторизации введённый пароль должен сверяться с хешем, хранящимся в базе данных.

Разграничение прав доступа.  
В системе предусмотрены роли заказчика и исполнителя. После успешной авторизации пользователь должен получать доступ только к тем функциям, которые соответствуют его роли. Заказчик должен иметь возможность создавать бронирования, работать с избранным, вести переписку и оставлять отзывы об исполнителях. Исполнитель должен иметь возможность управлять услугами, расписанием, входящими заявками и оставлять отзывы о заказчиках. Проверка прав доступа должна выполняться на серверной стороне перед выполнением любого действия.

Защита пользовательских данных.  
Пользователь не должен иметь возможности просматривать, изменять или удалять данные, не относящиеся к его учётной записи. Это относится к профилям, заявкам, сообщениям, отзывам, услугам и иным объектам системы.

Защита от SQL-инъекций.  
Все запросы к базе данных должны выполняться с использованием параметризованных запросов и подготовленных выражений. Это исключает возможность внедрения произвольного SQL-кода через пользовательский ввод.

Безопасность загрузки файлов.  
Файлы, загружаемые пользователями в систему, должны проходить проверку типа, размера и допустимости использования. Система должна ограничивать загрузку недопустимых форматов и предотвращать сохранение потенциально опасных файлов.

Защита сеанса работы.  
Система должна обеспечивать корректную работу пользовательских сессий, исключающую несанкционированное использование функций авторизованного пользователя после завершения сеанса или выхода из системы.

\subsection{Требования к программной и технической совместимости}

Система должна функционировать в стандартной web-среде и не предъявлять чрезмерных требований к программному и аппаратному обеспечению.

Операционная система.  
Серверная часть системы может функционировать под управлением операционных систем Windows или Linux при наличии установленного web-сервера Apache, интерпретатора PHP и системы управления базами данных MySQL. Клиентская часть системы не зависит от конкретной операционной системы и может использоваться на устройствах под управлением Windows, Linux, macOS, Android и iOS.

Аппаратные требования.  
Для работы серверной части системы требуется персональный компьютер или сервер, обеспечивающий:
\begin{itemize}
	\item процессор с тактовой частотой не ниже 2 ГГц;
	\item оперативную память объёмом не менее 4 ГБ;
	\item наличие свободного дискового пространства для размещения программных файлов, базы данных и загружаемых пользователями изображений;
	\item постоянный доступ к сети Интернет или локальной сети в зависимости от режима эксплуатации.
\end{itemize}

Программное обеспечение.  
Для реализации программной системы должны использоваться следующие программные средства:
\begin{itemize}
	\item язык разметки HTML;
	\item каскадные таблицы стилей CSS;
	\item язык программирования JavaScript;
	\item язык программирования PHP;
	\item система управления базами данных MySQL;
	\item web-сервер Apache.
\end{itemize}

Для работы клиентской части системы требуется любой современный web-браузер с поддержкой HTML5, CSS3 и JavaScript, например Google Chrome, Mozilla Firefox, Microsoft Edge, Safari или иной совместимый браузер.

Сетевые требования.  
Система должна обеспечивать возможность работы через web-интерфейс по протоколу HTTP или HTTPS. При размещении на сервере или в локальной сети пользователи должны получать доступ к приложению через браузер без необходимости установки дополнительного клиентского программного обеспечения.

Требования к размещению и сопровождению.  
Система должна поставляться в виде набора файлов web-приложения и базы данных, пригодных для развёртывания на сервере с поддержкой PHP и MySQL. Для сопровождения системы должны быть доступны исходные файлы проекта, база данных, загружаемые пользовательские материалы и резервные копии данных.

\subsection{Требования к оформлению документации}

Разработка программной документации и программного изделия должна производиться согласно ГОСТ 19.102--77 и ГОСТ 34.601--90. Единая система программной документации.